% TODO (author): fill author block, affiliations, emails, and CRediT statement before submission.
\documentclass[conference]{IEEEtran}
\IEEEoverridecommandlockouts
\usepackage{cite}
\usepackage{amsmath,amssymb,amsfonts}
\usepackage{graphicx}
\usepackage{textcomp}
\usepackage{xcolor}
\usepackage{booktabs}
\usepackage{url}
\def\BibTeX{{\rm B\kern-.05em{\sc i\kern-.025em b}\kern-.08em T\kern-.1667em\lower.7ex\hbox{E}\kern-.125em X}}
\begin{document}

\title{An Interpretable, Fairness-Aware Machine Learning Framework for Diabetes Screening from Population Health Surveys: Leakage, Calibration, and Fairness Audits}

\author{\IEEEauthorblockN{[Author Name]}
\IEEEauthorblockA{\textit{[Department]} \\
\textit{[Institution]} \\
[City], [Country] \\
[email]}}

\maketitle

\begin{abstract}
Diabetes affects hundreds of millions of people worldwide, yet a large share of cases remain undiagnosed, and machine learning models fitted to public health survey data carry failure modes that are invisible in headline accuracy. This paper audits an interpretable screening framework for diabetes on the 2015 Behavioral Risk Factor Surveillance System (BRFSS; $n = 253{,}680$, 21 features) and reports what the audits change. Logistic regression, random forest, and histogram-based gradient boosting are compared over five seeds; discrimination, calibration, and subgroup fairness are assessed; and three protocols---plain split, deduplicated test, and grouped held-out---are compared to bound duplicate-profile leakage. Leakage is small but not nil: the deduplicated test protocol lowers gradient-boosting AUC from $0.8269$ to $0.8177$, and 1{,}963 rows share a covariate vector with a contradicting label. Class-balancing weights leave the model badly miscalibrated (ECE $0.2367$, mean predicted probability $0.376$ against prevalence $0.139$); Beta, isotonic, and Platt recalibration reduce ECE to $0.0026$--$0.0083$, and the 2015 calibration map transports to 2017--2021 with an ECE penalty below $0.006$. Fairness gaps are quantified with bootstrap intervals and group supports: the sex TPR gap is small ($0.0241$, CI $[0.0066, 0.0417]$), whereas the age gap of $0.7889$ falls to $0.1436$ once bands with fewer than 500 positives are excluded. Group-specific thresholds targeting 0.85 sensitivity reduce the sex TPR gap to $0.0014$ and the education gap from $0.2400$ to $0.1699$, at a specificity cost. We frame the tool as a screening-level instrument that complements, rather than replaces, clinical assessment, and we state precisely where the data cannot support population-level claims.
\end{abstract}

\begin{IEEEkeywords}
diabetes screening, machine learning, calibration, fairness auditing, SHAP, BRFSS, health equity, ethical AI
\end{IEEEkeywords}

\section{Introduction}

\subsection{Problem context and practical requirements}
Diabetes is among the most prevalent chronic diseases globally, and a substantial fraction of cases are undiagnosed. Population screening can surface at-risk individuals, but any such tool must be reproducible on open data, transparent in how predictions are formed, honest about what it does not capture, and explicitly bounded against use for individual clinical diagnosis. This paper audits a survey-based screening tool against these requirements.

\subsection{Determinants of diabetes}
Beyond age and body mass index, diabetes risk reflects blood pressure, cholesterol, physical activity, diet, alcohol use, mental health, and socioeconomic conditions. Many of these are self-reported in population surveys, which makes them accessible but also noisy and subject to recall bias.

\subsection{Methodological approaches}
Prior work has applied logistic regression and, more recently, machine learning classifiers to the BRFSS, with AUC values roughly between 0.72 and 0.80 \cite{xie2019}. Related work has expanded the evidence base in three directions: non-invasive dysglycemia screening on NHANES with calibration and decision-curve analysis \cite{sun2026beyond}; large-scale external validation of NHANES-trained diabetes models on BRFSS, documenting a $9.7\%$ relative AUC drop under distribution shift and marked age-related disparities \cite{pall2026comprehensive}; and parsimonious feature selection on BRFSS, reducing 21 features to 10 with a copula tail-dependence filter at AUC $0.823$ \cite{aich2026copula}.

\subsection{Why audits change the conclusions}
Reported accuracy is a weak summary of a screening model, and several widely used claims about survey-based diabetes models are not reproducible under a leakage-audited protocol. On this same 253{,}680-row dataset, published pipelines report AUC near $0.99$ \cite{pall2026comprehensive}; we measure $0.827$ and show that duplicate profiles, class-weight-induced miscalibration, and support-starved fairness strata each move a headline number materially. The purpose of this paper is therefore to audit rather than to chase AUC: we report discrimination together with calibration, fairness with inference and supports, and we quantify each threat to the headline metric.

\subsection{Research questions and contributions}
We address three questions: (1) \textbf{Driver identification:} which features most strongly associate with diabetes status, and how stable are these attributions across seeds and survey waves? (2) \textbf{Fairness:} how do performance and error rates differ across sex, age, income, and education groups, and how much of each gap is support starvation? (3) \textbf{Deployment utility:} what operating point, and what calibration, does a screening program actually need? Our contributions are: (i) a reproducible multi-seed pipeline on a large public dataset; (ii) a three-protocol leakage audit with a reconciled duplicate definition; (iii) a calibration audit across four calibrators with per-wave transport analysis; (iv) fairness quantification with supports, bootstrap inference, design-effect sensitivity, and a support-restricted age estimate; (v) an implemented fairness mitigation through group-specific thresholds; (vi) SHAP explanations with stability metrics and comparison against alternative selectors; (vii) out-of-time validation on 2017--2021 waves; and (viii) decision-curve net benefit. The tool is positioned as screening-level and complementing, not replacing, clinical assessment.

\section{Materials and Methods}

\subsection{Data source and scope}
We use the CDC Diabetes Health Indicators dataset \cite{cdc2017}, derived from the 2015 BRFSS \cite{brfss}, a telephone survey conducted annually by the CDC since 1984. The cleaned release contains 253{,}680 listwise-complete respondents and 21 features; license CC0 (public domain), with raw BRFSS data in the US Government public domain. The target is \texttt{Diabetes\_binary} (0 = no diabetes, 1 = prediabetes or diabetes; prevalence 13.93\%). Two features of the release constrain our claims. First, race/ethnicity is absent, so fairness analysis is limited to sex, age, income, and education; race-level analysis requires the raw BRFSS files. Second, the release carries no survey design variables (weights, strata, PSUs), so population-level claims are not supported and we quantify only the sensitivity of our uncertainty to assumed design effects.

\subsection{Variables}
Table~\ref{tab:vars} groups the 21 features. Several are proxies: general health status stands in for clinical status, and income and education stand in for socioeconomic position. We treat these as proxies throughout and say so where results depend on them.

\begin{table}[h]
\caption{Feature groups.}
\label{tab:vars}
\centering
\begin{tabular}{ll}
\toprule
Group & Features \\
\midrule
Clinical & HighBP, HighChol, CholCheck, Stroke, HeartDiseaseorAttack, GenHlth, BMI \\
Lifestyle & Smoker, PhysActivity, Fruits, Veggies, HvyAlcoholConsump \\
Access & AnyHealthcare, NoDocbcCost \\
Mental/physical & MentHlth, PhysHlth, DiffWalk \\
Demographic & Sex, Age, Education, Income \\
\bottomrule
\end{tabular}
\end{table}

\subsection{Data preparation and sample flow}
Raw: 253{,}680 respondents $\rightarrow$ no exclusions (no missing values) $\rightarrow$ 253{,}680 final analytical sample. The \texttt{ID} column was dropped; no row was removed. Stratified 80/20 train/test split (seed 42) yields 202{,}944 train and 50{,}736 test cases, with a 64/16/20 train/calibration/test split for calibration work. All preprocessing is fit on the training portion only.

\subsection{Feature engineering}
Features are used as provided (binary or ordinal integers). Continuous-valued fields (BMI, MentHlth, PhysHlth, GenHlth) are standardized for logistic regression only; tree models use raw values. No derived features are added, keeping the screening questionnaire short.

\subsection{Validation strategy}
Stratified 5-fold cross-validation on the training portion for model development; a held-out stratified test set for final reporting. Variance is reported across five independent seeds (full re-split and re-train) as mean $\pm$ std. We acknowledge the single-split limitation.

\subsection{Model selection}
Three models: (M0) logistic regression with L2 penalty (interpretable baseline), (M1) random forest (50--100 trees, max depth 20), (M2) scikit-learn's histogram-based gradient boosting (150 iterations, 31 leaves). Class weights are balanced to address the 13.93\% prevalence. The humble-baseline move: logistic regression is included deliberately and reported honestly.

\subsection{Hyperparameter optimization}
Baseline settings are fixed and light. To verify that these do not disadvantage any baseline, we additionally run a randomized search (20 iterations for gradient boosting, 6 for random forest, 10 for logistic regression; 3-fold cross-validation on the training portion only) and report the best configuration found.

\subsection{Evaluation framework}
Discrimination: ROC-AUC, average precision (AUPRC), F1, accuracy. Calibration: Brier score, expected calibration error (ECE, 10 equal-width bins), and maximum calibration error (MCE), overall and per group, before and after recalibration with isotonic, Platt (logistic on $p$), Beta, and Venn--Abers calibrators. Leakage control: three held-out protocols (plain stratified split; deduplicated test set excluding any row whose covariate vector occurs in training; grouped held-out assignment with identical covariate vectors confined to one fold) plus a full deduplicated refit. Fairness: TPR gap, FPR gap, equalized-odds difference \cite{hardt2016}, and selection-rate gap, each with per-group supports and 300-resample bootstrap confidence intervals, reported with sensitivity to assumed design effects. Model comparison: paired bootstrap of the test-set AUC difference (2{,}000 resamples). Utility: decision-curve net benefit across thresholds. Ablations by feature group and by training-only feature ranking.

\subsection{Explainability method}
TreeSHAP \cite{shap2017} for the gradient boosting model on a 2{,}000-row test sample; mean absolute SHAP per feature; top-10 stability reported as mean pairwise Jaccard overlap across five seeds and across the 2017--2021 waves. Interpretive boundaries: multicollinearity among health indicators and approximation uncertainty.

\subsection{Scenario probe design}
Bounded perturbation probes: each modifiable feature is shifted to a specified value, clipped to its historical range, and the induced change in predicted probability is reported with a paired bootstrap confidence interval, both on raw and on recalibrated scores and broken down by subgroup. These probes characterise model behaviour and are not projections of causal effects.

\subsection{Software implementation and reproducibility}
Python 3.11.8; scikit-learn 1.5.2 \cite{sklearn}; XGBoost 2.1.3 \cite{xgboost}; SHAP 0.46.0; pandas 2.2.3; NumPy 2.1.3; matplotlib 3.9.2. Pinned versions are in \texttt{requirements.txt}; seeds are fixed (42, 7, 123, 2024, 99). All run manifests, per-seed logs, and a SHA-256 hash manifest of every artifact are released with the code \cite{repo}.

\section{Results}

\subsection{Predictive performance and model comparison}
Table~\ref{tab:perf} reports test performance over five seeds with AUPRC added, since screening operates under class imbalance. Gradient boosting leads on both metrics (ROC-AUC $0.8288 \pm 0.0014$, AUPRC $0.4229$), and logistic regression trails by less than $0.008$ AUC. A paired bootstrap of the test-set AUC difference (2{,}000 resamples) shows the separations are statistically significant yet practically negligible: HGB exceeds random forest by $\Delta$AUC $0.0034$ (95\% CI $[0.0023, 0.0044]$) and logistic regression by $0.0073$ (CI $[0.0058, 0.0088]$), all at $p<0.001$. A randomized search on the training portion returns best CV AUC $0.8303$ (HGB), $0.8275$ (RF), and $0.8233$ (LR), preserving the ordering, so the fixed settings did not disadvantage any baseline. We read this as the quantitative basis for the humble-baseline position: the differences are far below the $0.02$--$0.03$ margin that would change a screening decision, so an interpretable logistic baseline remains a defensible deployment option.

\begin{table}[h]
\caption{Test performance (mean $\pm$ std, 5 seeds) and paired AUC comparison.}
\label{tab:perf}
\centering
\begin{tabular}{lcccc}
\toprule
Model & Accuracy & F1 & ROC-AUC & AUPRC \\
\midrule
Logistic Regression & $0.7308 \pm 0.0013$ & $0.4418 \pm 0.0015$ & $0.8212 \pm 0.0014$ & 0.3926 \\
Random Forest & $0.7458 \pm 0.0008$ & $0.4515 \pm 0.0007$ & $0.8261 \pm 0.0015$ & 0.4153 \\
Hist. Grad. Boosting & $0.7203 \pm 0.0025$ & $0.4426 \pm 0.0023$ & $0.8288 \pm 0.0014$ & 0.4229 \\
\midrule
\multicolumn{5}{@{}l}{\textit{Paired bootstrap, HGB minus alternative (95\% CI), $p<0.001$ in both cases}} \\
HGB $-$ RF & --- & --- & $+0.0034\ (0.0023, 0.0044)$ & --- \\
HGB $-$ LR & --- & --- & $+0.0073\ (0.0058, 0.0088)$ & --- \\
\bottomrule
\end{tabular}
\end{table}

\subsection{Leakage audit and duplicate reconciliation}
We reconcile two duplicate counts that are frequently conflated. Under a strict definition---rows identical in all columns, including the outcome---24{,}206 rows (9.54\%) repeat an earlier row. Under a covariate-only definition---rows whose 21-feature vector repeats an earlier row regardless of outcome---25{,}772 rows (10.16\%) do so, leaving 227{,}908 distinct covariate vectors. The difference matters: 1{,}963 rows carry a covariate vector identical to an earlier row but a \emph{contradicting} label, which is a direct measure of label inconsistency in the release and is reported here for the first time. Table~\ref{tab:leak} compares three held-out protocols. The deduplicated test protocol (test rows with no covariate-vector overlap with training, $n = 43{,}900$) lowers AUC by $0.009$--$0.010$ for every model (HGB $0.8269 \rightarrow 0.8177$), while the grouped held-out protocol, which uses all 253{,}680 rows for evaluation, yields $0.8301$. The grouped-versus-plain 5-fold comparison differs by only $0.0001$. We therefore conclude that duplicate profiles do not explain published AUCs near $0.99$ on this dataset, but that the leakage-free estimate for our headline model is $0.8177$, not $0.8269$, and we report the latter only with the protocol attached.

\begin{table}[h]
\caption{Three held-out protocols (HGB unless noted).}
\label{tab:leak}
\centering
\begin{tabular}{lccc}
\toprule
Protocol & Test $n$ & ROC-AUC & AUPRC \\
\midrule
P1: plain stratified 80/20 --- LR & 50{,}736 & 0.8196 & 0.3926 \\
P1: plain stratified 80/20 --- RF & 50{,}736 & 0.8236 & 0.4205 \\
P1: plain stratified 80/20 --- HGB & 50{,}736 & 0.8269 & 0.4231 \\
P2: deduplicated test --- LR & 43{,}900 & 0.8096 & 0.3994 \\
P2: deduplicated test --- RF & 43{,}900 & 0.8143 & 0.4284 \\
P2: deduplicated test --- HGB & 43{,}900 & 0.8177 & 0.4309 \\
P3: grouped held-out (5-fold) --- HGB & 253{,}680 & 0.8301 & 0.4326 \\
\bottomrule
\end{tabular}
\end{table}

\subsection{Calibration audit and post-hoc recalibration}
Discrimination alone is insufficient for screening, where a communicated probability must correspond to an observed frequency. Because class-balancing weights distort the probability scale, the uncalibrated model is badly miscalibrated (Brier $0.1757$, ECE $0.2367$, MCE $0.4201$; mean predicted probability $0.376$ against observed prevalence $0.139$). Table~\ref{tab:cal} compares four calibrators fitted on a held-out calibration split. Beta calibration is best on all three metrics (Brier $0.0974$, ECE $0.0026$, MCE $0.0508$), closely followed by isotonic (ECE $0.0031$ but MCE $0.2500$) and Platt (ECE $0.0083$, MCE $0.0559$); Venn--Abers performs worst here (ECE $0.1534$, mean predicted probability $0.2815$), because the self-consistent prior correction is poorly matched to a strongly over-confident boosted ensemble. All are monotone, so AUC is unchanged ($0.8265$). Per-group maximum ECE falls from $0.2505$ to $0.0045$ (sex), $0.3403$ to $0.0208$ (age), $0.3093$ to $0.0343$ (income), and $0.3478$ to $0.0342$ (education) under isotonic recalibration. We regard recalibration as a practical precondition rather than a refinement: a balanced-weight model must be recalibrated before any probability is communicated to a user or clinician.

\begin{table}[h]
\caption{Calibration before and after post-hoc recalibration (HGB, 2015 test, $n=50{,}736$).}
\label{tab:cal}
\centering
\begin{tabular}{lcccc}
\toprule
Calibrator & Brier & ECE & MCE & Mean pred. \\
\midrule
Uncalibrated & 0.1757 & 0.2367 & 0.4201 & 0.376 \\
Isotonic & 0.0976 & 0.0031 & 0.2500 & 0.140 \\
Platt (logistic) & 0.0977 & 0.0083 & 0.0559 & 0.139 \\
Beta & \textbf{0.0974} & \textbf{0.0026} & \textbf{0.0508} & 0.139 \\
Venn--Abers & 0.1266 & 0.1534 & 0.2420 & 0.282 \\
\midrule
Observed prevalence & --- & --- & --- & 0.139 \\
\bottomrule
\end{tabular}
\end{table}

\begin{figure}[h]
\centering
\includegraphics[width=\linewidth]{calibration_curve.png}
\caption{Reliability diagrams before and after isotonic recalibration (HGB, 2015 test).}
\label{fig:cal}
\end{figure}

\subsection{Out-of-time validation and calibration transport}
Training on the 2015 wave and testing on later waves (Table~\ref{tab:temp}) yields AUC $0.8338$ (2017), $0.8256$ (2019), and $0.8198$ (2021) with a mild downward trend. We also asked whether the 2015 calibration map can be transported rather than refit per wave, since refitting per wave is not possible in a prospective deployment. Applying the 2015-fitted isotonic map to later waves gives ECE $0.0064$, $0.0062$, and $0.0097$ against raw values of $0.2240$, $0.2326$, and $0.2236$; refitting the calibrator on each wave improves ECE by only $0.0035$, $0.0030$, and $0.0056$ (Brier penalties below $0.002$), so the transport penalty is small. Maximum subgroup ECE also remains low after transport (age $0.0168$--$0.0246$). The threshold required to hold 85\% sensitivity moves only from $0.420$ to $0.395$ across waves, so neither the operating point nor the calibration map requires per-wave re-estimation. For context, the $0.014$ AUC drop from 2017 to 2021 is milder than the $9.7\%$ relative drop reported when an NHANES-trained model was validated on BRFSS \cite{pall2026comprehensive}, which is consistent with the absence of a cross-survey shift.

\begin{table}[h]
\caption{Out-of-time validation with transported (2015-fitted) calibration.}
\label{tab:temp}
\centering
\begin{tabular}{lcccccc}
\toprule
Test wave & $n$ & Prev. & ROC-AUC & Raw ECE & Transport ECE & Thr. (sens.\ 0.85) \\
\midrule
2015 & 50{,}736 & 0.139 & 0.8265 & 0.2367 & 0.0031 & --- \\
2017 & 275{,}854 & 0.141 & 0.8338 & 0.2240 & 0.0064 & 0.420 \\
2019 & 243{,}317 & 0.145 & 0.8256 & 0.2326 & 0.0062 & 0.420 \\
2021 & 236{,}378 & 0.142 & 0.8198 & 0.2236 & 0.0097 & 0.395 \\
\bottomrule
\end{tabular}
\end{table}

\subsection{Fairness with supports, inference, and design-effect sensitivity}
We quantify subgroup disparities with per-group supports and bootstrap confidence intervals (300 resamples) rather than point estimates alone, following recent audit practice for clinical risk models \cite{roy2026kaisen}. Table~\ref{tab:fair} reports the TPR gap, the equalized-odds FPR gap, and the smallest positive count in any group. We caution against over-interpreting gap metrics without base-rate context \cite{mehrabi2021}. The sex TPR gap is small ($0.0241$, 95\% CI $[0.0066, 0.0417]$); the age TPR gap is very large ($0.7889$, CI $[0.6421, 0.8774]$), but the smallest age band contains only 14 positives.

Because the cleaned release carries no design variables, we cannot estimate the true design effect, so we report how the subgroup uncertainty would widen if the design effect took typical values for a stratified multistage survey (Table~\ref{tab:fair}). The widest TPR-gap interval grows from $0.235$ at DEFF $=1$ to $0.407$ at DEFF $=3$; the qualitative conclusions---small sex gap, large age gap---are unchanged, but we note that our intervals are, if anything, optimistic.

\begin{table}[h]
\caption{Subgroup disparities with support and bootstrap inference (HGB, threshold 0.5).}
\label{tab:fair}
\centering
\begin{tabular}{lccccc}
\toprule
Stratifier & TPR gap (95\% CI) & FPR gap & Min $n$ & Min $n_{pos}$ & CI width, DEFF 3 \\
\midrule
Sex & 0.0241 (0.0066, 0.0417) & 0.0643 & 22{,}205 & 3{,}384 & 0.074 \\
Age (13 bands) & 0.7889 (0.6421, 0.8774) & 0.5320 & 1{,}207 & 14 & 0.407 \\
Income (8 levels) & 0.2849 (0.2571, 0.3173) & 0.4229 & 1{,}941 & 475 & 0.268 \\
Education (6 levels) & 0.2436 (0.2107, 0.2781) & 0.4738 & 795 & 221 & 0.330 \\
\bottomrule
\end{tabular}
\end{table}

\subsection{Support-restricted age disparity and severity stratification}
Repeating the age computation under increasing minimum-support thresholds (Table~\ref{tab:age}) shows the gap falls from $0.7879$ over all 13 bands to $0.1436$ when only bands with at least 500 positives are retained, while the widest within-group 95\% confidence interval for TPR narrows from $0.340$ to $0.072$. We therefore report $0.1436$ (seven bands, $\ge 500$ positives) as the defensible estimate and treat the headline $0.7889$ as support-starved; the residual gap is nevertheless real, so a single global threshold is not age-equitable. Because the release collapses prediabetes and diabetes into one label, we cannot separate the two conditions, and we instead test whether the model discriminates and calibrates within severity strata defined by available proxies (Table~\ref{tab:age}). Within-stratum AUC ranges from $0.7199$ to $0.8263$ and post-recalibration ECE from $0.0025$ to $0.0134$, so the model neither collapses nor inflates its discrimination in the sicker strata, and recalibration holds within strata.

\begin{table}[h]
\caption{Age TPR gap under minimum support, and severity-stratified performance (HGB).}
\label{tab:age}
\centering
\begin{tabular}{lccc|lccc}
\toprule
\multicolumn{4}{c|}{\textit{Age gap vs.\ minimum support}} & \multicolumn{3}{c}{\textit{Severity strata}} \\
Min.\ pos. & Bands & TPR gap & Widest CI & Stratum & $n$ & Prev. & AUC \\
\midrule
1 & 13 & 0.7879 & 0.340 & GenHlth 1--2 & 26{,}988 & 0.057 & 0.8145 \\
20 & 12 & 0.5426 & 0.340 & GenHlth 3 & 15{,}064 & 0.177 & 0.7443 \\
50 & 11 & 0.5426 & 0.229 & GenHlth 4--5 & 8{,}684 & 0.330 & 0.7199 \\
100 & 10 & 0.4206 & 0.154 & BMI $<30$ & 33{,}123 & 0.092 & 0.8263 \\
500 & 7 & \textbf{0.1436} & 0.072 & BMI $\ge30$ & 17{,}613 & 0.227 & 0.7802 \\
\bottomrule
\end{tabular}
\end{table}

\subsection{Operating points, constrained selection, and a fairness mitigation}
For screening, sensitivity matters more than overall accuracy. We first characterise the trade-off directly: a single global threshold selected to maximise specificity subject to overall sensitivity $\ge 0.85$ and a bound $\delta$ on the worst-group TPR gap (Table~\ref{tab:thr}). At the default threshold $0.50$ the operating point gives sensitivity $0.799$, specificity $0.703$, and a worst-group TPR gap of $0.545$. Binding the gap to $\delta = 0.30$ forces the threshold down to $0.140$, raising sensitivity to $0.975$ but collapsing specificity to $0.344$; even $\delta = 0.50$ leaves specificity at only $0.621$. A single global threshold is therefore a poor fairness instrument for this model, which is the substantive reason we next implement group-specific thresholds.

Group-specific thresholds, fitted on the calibration split to deliver 0.85 sensitivity within each sex, age, income, and education stratum, reduce the sex TPR gap from $0.0228$ to $0.0014$, the education gap from $0.2400$ to $0.1699$, and the income gap from $0.2862$ to $0.2300$, with the age gap nearly unchanged ($0.5426 \rightarrow 0.5263$). Overall sensitivity rises from $0.795$ to $0.857$ and specificity falls from $0.708$ to $0.637$, and the age FPR gap widens from $0.5198$ to $0.6341$. This is an implemented mitigation rather than an audit: it closes the sex gap almost completely, and it prices the remaining trade-off explicitly. It is not a general solution, since it cannot recover the age gap without accepting a far more permissive threshold.

\begin{table}[h]
\caption{Operating points: global, constrained, and group-specific thresholds.}
\label{tab:thr}
\centering
\begin{tabular}{lcccc}
\toprule
Policy & Threshold & Sens. & Spec. & Worst TPR gap \\
\midrule
Global, default & 0.500 & 0.799 & 0.703 & 0.545 \\
Constrained, $\delta=0.50$ & 0.400 & 0.867 & 0.621 & 0.494 \\
Constrained, $\delta=0.30$ & 0.140 & 0.975 & 0.344 & 0.278 \\
Constrained, $\delta=0.10$ & 0.045 & 0.997 & 0.147 & 0.069 \\
\midrule
\multicolumn{5}{@{}l}{\textit{Group-specific thresholds, per-stratum target sensitivity 0.85}} \\
\quad Sex & per group & --- & --- & 0.0014 \\
\quad Income & per group & --- & --- & 0.2300 \\
\quad Education & per group & --- & --- & 0.1699 \\
\quad Age & per group & --- & --- & 0.5263 \\
\quad All (overall) & per group & 0.857 & 0.637 & --- \\
\bottomrule
\end{tabular}
\end{table}

\subsection{Explainability findings and stability}
SHAP (Fig.~\ref{fig:shap}) ranks general health status (0.638), high blood pressure (0.541), age (0.416), BMI (0.382), and high cholesterol (0.313) as the leading contributors, followed by income (0.136), sex (0.125), cholesterol check (0.092), mental health (0.063), heavy alcohol use (0.060), heart disease (0.059), difficulty walking (0.056), and education (0.045). These align with established risk factors and support the model's face validity, although several are consequences rather than antecedents of diabetes, which limits their use for early-detection reasoning. The top-10 ordering is stable: mean pairwise Jaccard overlap of the top-10 sets is $0.964$ across the five seeds (minimum $0.818$, differing only in the tenth position) and $0.818$ between 2015 and each later wave, where the substitution is heavy alcohol use for difficulty walking. We therefore report the feature ranking as a property of the data period rather than of the seed.

\begin{figure}[h]
\centering
\includegraphics[width=0.85\linewidth]{shap_top_features.png}
\caption{Top 12 features by mean $|$SHAP$|$ (HGB, 2{,}000-row test sample).}
\label{fig:shap}
\end{figure}

\subsection{Feature subset analyses}
Table~\ref{tab:abl} combines the feature-group ablations with the minimal-questionnaire experiment and a comparison of ranking procedures. Clinical-history features carry most of the signal (AUC $0.8241$ versus $0.7104$ for demographics alone), while lifestyle features alone are weaker ($0.7475$). Using training-only SHAP rankings, five features retain AUC $0.8186$ against $0.8269$ for all 21, and eleven reach $0.8257$. Comparing three training-only ranking procedures at $k=10$ against a random-subset baseline, SHAP attains $0.8248$, correlation filtering $0.8218$, mutual information $0.8213$, and ten random features $0.7784 \pm 0.0108$, so SHAP is strongest but only marginally over the other filters. For context, a copula tail-dependence filter on this same dataset has been reported at $0.823$ with 10 features \cite{aich2026copula}, consistent with our range. Because several top features proxy the same constructs, we do not claim that a ten-item form is clinically validated.

\begin{table}[h]
\caption{Feature subset analyses (HGB, 2015 test).}
\label{tab:abl}
\centering
\begin{tabular}{lccc|lccc}
\toprule
\multicolumn{4}{c|}{\textit{Feature-group ablation}} & \multicolumn{3}{c}{\textit{Subset and selector comparison}} \\
Feature group & $n$ & ROC-AUC & F1 & Setting & $n$ & ROC-AUC & F1 \\
\midrule
Demographic only & 4 & 0.7104 & --- & Top-5 SHAP & 5 & 0.8186 & 0.4305 \\
$+$ clinical history & 11 & 0.8241 & --- & Top-8 SHAP & 8 & 0.8231 & 0.4357 \\
$+$ lifestyle & 14 & 0.7475 & --- & Top-11 SHAP & 11 & 0.8257 & 0.4404 \\
Full & 21 & 0.8269 & 0.4411 & All features & 21 & 0.8269 & 0.4411 \\
 & & & & Mutual info $k{=}10$ & 10 & 0.8213 & --- \\
 & & & & Correlation $k{=}10$ & 10 & 0.8218 & --- \\
 & & & & Random $k{=}10$ & 10 & \multicolumn{2}{c}{$0.7784 \pm 0.0108$} \\
\bottomrule
\end{tabular}
\end{table}

\subsection{Decision-curve utility}
For a screening program the operating point must be justified by net benefit rather than by a generic accuracy target. Figure~\ref{fig:dca} gives decision-curve net benefit against threshold probability. The model dominates treat-all and treat-none across the clinically relevant range: at a threshold of $0.10$ net benefit is $0.068$ against $0.044$ for treat-all, and at $0.20$ it is $0.009$ against $-0.076$, with treat-all becoming actively harmful beyond $0.10$. The decision curve therefore supports the sensitivity-oriented operating points of Table~\ref{tab:thr} and indicates that the useful threshold range lies well below $0.5$, consistent with the strong over-prediction of the uncalibrated model.

\begin{figure}[h]
\centering
\includegraphics[width=0.8\linewidth]{decision_curve.png}
\caption{Decision-curve net benefit versus threshold probability (HGB, 2015 test).}
\label{fig:dca}
\end{figure}

\subsection{Scenario probes}
We performed bounded perturbation probes on the test set, clipping each modified feature to its historical range; these are sensitivity probes of model behaviour, not projections of causal effects. A $+5$ BMI shift raises mean predicted probability by $0.0694$ on raw scores but by only $0.0410$ on recalibrated scores, and a one-point worsening of general health raises it by $0.0985$ raw and $0.0499$ recalibrated: recalibration compresses the response, so probe magnitudes should be quoted on the scale of the scores actually deployed. The smoking and physical-activity probes are small and not distinguishable from zero (smoker $0\rightarrow1$: $+0.0002$, CI $[-0.0208, 0.0128]$; physical activity $1\rightarrow0$: $+0.0036$, CI $[-0.0042, 0.0230]$), and setting heavy alcohol use to present \emph{decreases} mean predicted probability by $0.1034$ (CI $[-0.3094, 0.0031]$), with only $4.4\%$ of individuals increasing. We report this pattern explicitly: the fitted model encodes the historical associations of this survey, including the non-monotone association with heavy alcohol use, and we do not read these probes as behavioural or clinical guidance.

\section{Discussion}

\subsection{What the audits changed}
The audit work materially altered five claims. First, the age TPR gap is not a single number: it is $0.7889$ over all 13 bands but $0.1436$ over the seven bands with at least 500 positives, so we report the support-restricted figure. Second, the model is not usable as a probability estimator without recalibration: class-balancing weights push mean predicted probability from $0.139$ to $0.376$, and Beta recalibration restores ECE to $0.0026$, with the 2015 map transporting to later waves at a penalty below $0.006$. Third, the AUC differences between the three models, though statistically significant, are too small to motivate a more complex model. Fourth, the headline AUC is protocol-dependent: the leakage-free estimate is $0.8177$, not $0.8269$, and published figures near $0.99$ on this dataset are not reproducible under any of our three protocols. Fifth, a single global threshold is a poor fairness instrument here---bounding the worst-group TPR gap to $0.10$ costs $0.56$ in specificity---whereas group-specific thresholds close the sex gap almost completely at a specificity cost of $0.07$.

\subsection{Contextual positioning within the literature}
Our discrimination (AUC $0.8288$, AUPRC $0.4229$) exceeds earlier BRFSS-based studies reporting 0.72--0.80 \cite{xie2019} and is close to NHANES-based non-invasive dysglycemia screening at AUC $0.820$ \cite{sun2026beyond}, but we now report AUPRC and decision-curve net benefit as well, since discrimination alone is not comparable across studies with different label definitions and prevalence. It is markedly more conservative than BRFSS studies reporting AUC near $0.99$ with stacked ensembles and resampling-based balancing \cite{pall2026comprehensive}; our leakage audit indicates that on this same 253{,}680-row dataset, duplicate profiles and fold construction move AUC by at most $0.010$, so we regard those figures as unlikely to be reproducible. Where an NHANES-trained model was externally validated on BRFSS, AUC fell from $0.794$ to $0.717$ under distribution shift \cite{pall2026comprehensive}; our in-domain 2015-to-later-wave transfer is milder ($0.8338$ to $0.8198$), consistent with the absence of a cross-survey shift. For feature reduction, a copula tail-dependence filter on this same dataset reports AUC $0.823$ with 10 features \cite{aich2026copula}, against our $0.8248$ for SHAP at $k=10$; direct comparison remains limited by differing years, features, splits, and label definitions, and we do not claim a decisive win.

\subsection{Practical deployment scenarios}
(1) Population-level risk stratification for outreach, using recalibrated probabilities and an explicit sensitivity target; (2) questionnaire triage in resource-limited settings, using the 8--10 item reduced form; (3) hypothesis generation for public-health research. Each requires corroboration with clinical measurements before any individual decision. Deployment should use group-specific rather than global thresholds, and should quote probabilities only after recalibration.

\subsection{Methodological limitations and validity boundaries}
Numbered: (1) self-reported cross-sectional data, no biomarkers (HbA1c, glucose), and several leading predictors are consequences rather than antecedents; (2) race/ethnicity absent from this release, limiting health-equity analysis and precluding race-stratified reliability checks; (3) no survey design weights, so population-level representativeness is not established, no design-based uncertainty is computed, and our intervals are optimistic by an unquantified design effect; (4) prediabetes and diabetes are collapsed into one label, and we cannot separate the two conditions, so the severity-stratified analysis is a proxy rather than a substitute; (5) the leakage-free estimate uses a deduplicated test set of $43{,}900$ rows, and residual support effects remain; (6) calibration is restored post hoc on a single 2015 split, and although the map transports well to 2017--2021, the 2023 wave lacks four features and was therefore excluded; (7) the randomized search budget is small and single-split; (8) group-specific thresholds were fitted on one calibration split and were not validated out of time; (9) association, not causation; (10) screening, not diagnosis. Design-aware modelling and design-based inference are the established remedy for population-level claims \cite{matabuena2024screening,wieczorek2023design}, and we identify them as prerequisites for any policy-facing use.

\subsection{Prospects for future development}
Incorporation of BRFSS design weights with design-based uncertainty quantification \cite{matabuena2024screening,wieczorek2023design}; inclusion of race/ethnicity from the raw BRFSS files with race-stratified reliability; extension of the out-of-time analysis to 2023 and to cross-survey validation against NHANES \cite{pall2026comprehensive}; per-group calibration within the group-specific-threshold policy, following recent audit recommendations \cite{roy2026kaisen}; and prospective evaluation with clinical biomarkers.

\section{Conclusions}
We audited an interpretable, fairness-aware screening framework for diabetes on 253{,}680 BRFSS respondents and report what the audits change. (1) Gradient boosting and random forest outperform logistic regression by a statistically significant but practically negligible margin ($\Delta$AUC $\le 0.0073$), so an interpretable baseline remains defensible. (2) The headline AUC is protocol-dependent: the leakage-free estimate is $0.8177$, and 1{,}963 rows carry a covariate vector with a contradicting label. (3) The balanced-weight model is badly miscalibrated (ECE $0.2367$) until recalibration (ECE $0.0026$ with Beta), which is a precondition for communicating any probability; the 2015 map transports to later waves at a penalty below $0.006$ in ECE. (4) The sex disparity is small, while the large age TPR gap is concentrated in sparse strata, falling from $0.7889$ to $0.1436$ under a minimum-support constraint, and group-specific thresholds close the sex gap to $0.0014$ at a specificity cost of $0.07$. (5) Decision-curve net benefit supports sensitivity-oriented operating points below $0.5$. Appropriate use context: the tool is a screening-level, hypothesis-generating instrument for population triage and research, not a validated diagnostic or causal instrument, and it does not yet support population-level prevalence claims.

\section*{Acknowledgments}
The authors thank the CDC for the BRFSS data and the UCI repository for hosting.

\section*{Data Availability}
The dataset is publicly available \cite{cdc2017}. All code, run manifests, logged results, and a SHA-256 hash manifest of every artifact are released at \cite{repo}.

\section*{Author Contributions}
[TBD: CRediT statement.]

\section*{Conflicts of Interest}
The authors declare no conflict of interest.

\begin{thebibliography}{99}

\bibitem{cdc2017}
CDC Diabetes Health Indicators [Dataset]. (2017). UCI Machine Learning Repository. \url{https://doi.org/10.24432/C53919} (accessed on 17 Oct. 2026).

\bibitem{xie2019}
Z. Xie, O. Nikolayeva, J. Luo, and D. Li, ``Building risk prediction models for type 2 diabetes using machine learning techniques,'' \textit{Prev. Chronic Dis.}, vol. 16, p. e130, 2019. \url{https://doi.org/10.5888/pcd16.190109}.

\bibitem{shap2017}
S. M. Lundberg and S.-I. Lee, ``A unified approach to interpreting model predictions,'' in \textit{Adv. Neural Inf. Process. Syst.}, vol. 30, pp. 4768--4777, 2017.

\bibitem{hardt2016}
M. Hardt, E. Price, and N. Srebro, ``Equality of opportunity in supervised learning,'' in \textit{Adv. Neural Inf. Process. Syst.}, vol. 29, pp. 3315--3323, 2016.

\bibitem{mehrabi2021}
N. Mehrabi, F. Morstatter, N. Saxena, K. Lerman, and A. Galstyan, ``A survey on bias and fairness in machine learning,'' \textit{ACM Comput. Surv.}, vol. 54, no. 6, art. 115, 2021. \url{https://doi.org/10.1145/3457607}.

\bibitem{sun2026beyond}
B. Sun, C. Zhang, K. Ji, and X. Lu, ``Beyond the blood draw: Explainable machine learning for non-invasive dysglycemia risk screening,'' \textit{arXiv:2606.16056}, 2026.

\bibitem{pall2026comprehensive}
R. S. Pall, S. Yadav, S. Bhalerao, S. Sahu, R. Ahluwalia, and B. Awadhiya, ``Comprehensive evaluation of machine learning for type 2 diabetes risk prediction: Large-scale external validation and fairness analysis,'' in \textit{Proc. IEEE CIPHER}, 2026. \url{https://doi.org/10.1109/CIPHER70417.2026.11523789}.

\bibitem{aich2026copula}
A. Aich, M. M. Murshed, S. Hewage, and A. Mayeaux, ``A copula based supervised filter for feature selection in machine learning driven diabetes risk prediction,'' \textit{Sci. Rep.}, 2026. \url{https://doi.org/10.1038/s41598-026-41874-9}.

\bibitem{matabuena2024screening}
M. Matabuena, J. C. Vidal, R. Ghosal, and J.-P. Onnela, ``Screening for diabetes mellitus in the U.S. population using neural network models and complex survey designs,'' \textit{arXiv:2403.19752}, 2024.

\bibitem{roy2026kaisen}
S. Roy, S. Girmachew, and N. Chavan, ``KAISEN: Reproducible subgroup fairness auditing for clinical risk models,'' \textit{arXiv:2607.28608}, 2026.

\bibitem{wieczorek2023design}
J. Wieczorek, ``Design-based conformal prediction,'' \textit{Survey Methodology}, vol. 49, no. 2, pp. 305--322, 2023.

\bibitem{brfss}
Centers for Disease Control and Prevention, ``Behavioral Risk Factor Surveillance System,'' 2015 annual data. \url{https://www.cdc.gov/brfss/annual_data/annual_2014.html} (accessed on 17 Oct. 2026).

\bibitem{sklearn}
F. Pedregola \textit{et al.}, ``Scikit-learn: machine learning in Python,'' \textit{J. Mach. Learn. Res.}, vol. 12, pp. 2825--2830, 2011.

\bibitem{xgboost}
T. Chen and C. Guestrin, ``XGBoost: a scalable tree boosting system,'' in \textit{Proc. ACM SIGKDD}, pp. 785--794, 2016.

\bibitem{repo}
A. Author, ``BRFSS diabetes screening: code, run manifests, and SHA-256 artifact manifest,'' GitHub repository. \url{https://github.com/Aldrin7/brfss-diabetes-screening-iccica2027} (accessed on 21 Oct. 2026).

\end{thebibliography}

\end{document}